% Evans 8.3: printed pp.454--458 / English PDF pp.471--475.
% Chinese: printed pp.320--323 / scan pp.338--341.
\Needspace{55mm}
\section{Regularity}
\label{sec:8-3}
\Needspace{45mm}
\subsection{Nonlinear difference quotients}
Consider a scalar weak solution $u\in H^1(U)$ of
\begin{equation}\label{eq:8-nonlinear-elliptic}
 -\operatorname{div}DL(Du)=f,\qquad f\in L^2(U),
\end{equation}
where $L\in C^2(\R^n)$ and
\begin{equation}\label{eq:8-uniform-hessian}
 \theta|\xi|^2\leq D^2L(p)\xi\cdot\xi\leq\Lambda|\xi|^2
 \quad(p,\xi\in\R^n),\qquad0<\theta\leq\Lambda.
\end{equation}
Subtracting the affine function $DL(0)\cdot p$ from $L$ leaves the weak
equation unchanged, because $\int_U DL(0)\cdot D\phi=0$ for
$\phi\in H^1_0$. We may therefore assume $DL(0)=0$, giving
$|DL(p)|\leq\Lambda|p|$.
\Needspace{55mm}
\begin{theorem}[Second derivatives; source \S8.3.1, Theorem 1]
\label{thm:8-H2}
For $V\Subset U$, $u\in H^2(V)$ and
\[
 \|D^2u\|_{L^2(V)}\leq C(\|f\|_{L^2(U)}+\|Du\|_{L^2(U)}).
\]
If $u\in H^1_0(U)$ and $\partial U$ is $C^2$, then
$u\in H^2(U)$ and $\|u\|_{H^2(U)}\leq C\|f\|_2$.
Constants depend on $\theta,\Lambda$, the domains and their indicated geometry.
\end{theorem}
\begin{expandedproof}[uniformly elliptic averaged Hessians]
Choose $V\Subset W\Subset U$ and a cutoff $\zeta$ supported in $W$,
equal to one on $V$. For a small step in direction $e_k$ let
$\delta_hv=(v(x+he_k)-v(x))/h$ and use the legal test
$\phi=-\delta_{-h}(\zeta^2\delta_hu)$.
It belongs to $H^1_0(U)$ because its support stays away from the boundary.
Discrete integration by parts gives
\[
 \int_U\delta_hDL(Du)\cdot D(\zeta^2\delta_hu)
 =-\int_U f\,\delta_{-h}(\zeta^2\delta_hu).
\]
The exact nonlinear linearization is
\[
 \delta_hDL(Du)=A_h\delta_hDu,\quad
 A_h=\int_0^1D^2L(Du+s(Du(x+he_k)-Du))\,ds.
\]
Hence $\theta I\leq A_h\leq\Lambda I$ independently of $h$.
Expanding the cutoff derivative, the principal integral bounds
$\theta\|\zeta\delta_hDu\|_2^2$. Its cross term costs
$\varepsilon\|\zeta\delta_hDu\|_2^2+C_\varepsilon\|\delta_hu\|_2^2$.
The forcing is controlled without differentiating $f$:
\[
 \begin{split}
 \left|\int f\,\delta_{-h}(\zeta^2\delta_hu)\right|
 &\leq C\|f\|_2\|D(\zeta^2\delta_hu)\|_2\\
 &\leq\varepsilon\|\zeta\delta_hDu\|_2^2+
 C_\varepsilon(\|f\|_2^2+\|Du\|_2^2).
 \end{split}
\]
Use $\|\delta_hu\|_2\leq\|Du\|_2$ on a slightly larger set and absorb
with $2\varepsilon<\theta$. The difference-quotient criterion from
\S5.8 then gives the interior estimate.
\end{expandedproof}
\Needspace{55mm}
\PDEAnnotation{平均 Hessian 保留椭圆性}{$A_h$ 是沿两点梯度连线的 Hessian 平均；一致上下界在平均后不变。不能只在某个未经控制的中间点套用局部凸性。}

\begin{expandedproof}[boundary and normal estimates]
This expands the boundary argument left to the reader in the source.
In a boundary chart write $x=\Phi(y)=(y',y_n+\gamma(y'))$, $v=u\circ\Phi$.
Its Jacobian is one and $v=0$ on $y_n=0$.
The transformed flux is
\[
 a(y,p)=D\Phi^{-1}\,DL(D\Phi^{-T}p),\qquad
 -\partial_{y_i}a_i(y,Dv)=\widetilde f.
\]
Bounded $D\Phi,D\Phi^{-1}$ preserve uniform monotonicity in $p$;
bounded $D^2\Phi$ gives $|D_ya(y,p)|\leq C|p|$, since $DL(0)=0$.
Tangential difference quotients now work on the half-chart. Their tests
have zero trace because $v$ and its tangential translates have zero trace.
Besides the averaged Hessian term, differencing $a$ produces a coefficient
error bounded by $C|Dv|$, which is absorbed as above. Thus all $v_{i\alpha}$,
$\alpha<n$, are in $L^2$, with bounds by $\|\widetilde f\|_2+\|Dv\|_2$.

To recover $v_{nn}$ without assuming it exists, let $w=a_n(y,Dv)$.
The known tangential derivatives and the flux chain rule give
$\partial_\alpha a_i(y,Dv)\in L^2$. The equation gives
$\partial_n w=-\widetilde f-\sum_{\alpha<n}\partial_\alpha a_\alpha\in L^2$.
For fixed $(y,p')$, $p_n\mapsto a_n(y,p',p_n)$ is strictly increasing
with derivative at least $c>0$ and at most $C$; it maps $\R$ onto $\R$.
Its inverse $p_n=\Psi(y,p',w)$ satisfies
\[
 |\Psi_w|+|\Psi_{p'}|\leq C,\qquad
 |\Psi_y|\leq C(|p'|+|\Psi|).
\]
Apply normal difference quotients to
$v_n=\Psi(y,D'v,w)$. The $L^2$ bounds for the normal derivatives of
$D'v$ and $w$ give an $L^2$ bound for the normal difference quotients
of $v_n$. The criterion proves $v_{nn}\in L^2$; this step avoids a circular
classical expansion of an equation before $H^2$ has been established.
Cutoffs, a finite chart cover and the interior estimate yield
$\|u\|_{H^2}\leq C(\|f\|_2+\|u\|_{H^1})$.
Finally testing the original equation by $u$, monotonicity and Poincar\'e give
\[
 \theta\|Du\|_2^2\leq\int_U DL(Du)\cdot Du=(f,u)
 \leq C\|f\|_2\|Du\|_2.
\]
This removes the lower-order norm and proves the global estimate.
\end{expandedproof}
\Needspace{55mm}
\PDEAnnotation{法向导数的非循环恢复}{先证明切向二阶导数，再由方程得到法向通量的弱导数，最后利用 $p_n\mapsto a_n$ 的一致单调可逆性恢复 $u_{nn}$。不能在已知 $u\in H^1$ 时直接展开二阶链式法则。}

\Needspace{45mm}
\subsection{Why further smoothness needs a new theorem}
For $f=0$, interior $H^2$ and the Lipschitz chain rule imply that
$v_k=u_{x_k}$ weakly solves
\[
 -\partial_i(A^{ij}(x)\partial_jv_k)=0,\qquad
 A^{ij}(x)=L_{p_ip_j}(Du(x)).
\]
These coefficients are only bounded measurable at this stage. The smooth
coefficient estimates of Chapter 6 cannot be applied to them.
The scalar De Giorgi--Nash--Moser theorem, used here as an external
regularity result, states that $H^1_{\rm loc}$ solutions of a uniformly
elliptic divergence equation with bounded measurable coefficients are
locally H\"older continuous. It gives $Du\in C^{0,\alpha}_{\rm loc}$.
For smooth $L$, $A$ is then H\"older. The external Schauder estimates
for linear divergence equations improve $v_k$ to $C^{1,\alpha}_{\rm loc}$;
iteration improves coefficient and solution regularity together, yielding
$u\in C^\infty(U)$.
This is a conditional bootstrap using those two named theorems, not a
proof of either deep estimate. For systems the scalar H\"older theorem
is unavailable in general; smoothness cannot be asserted by the same argument.
\Needspace{55mm}
\PDEAnnotation{线性化不意味着光滑系数}{$A(x)=D^2L(Du(x))$ 的正则性由未知梯度决定。先得到梯度的 H\"older 连续性，才能启动 Schauder 提升；多次形式求导不能绕过这一步。}
